\section{Conclusiones y Comentarios}

A partir del desarrollo de las presentes Bases de Cálculo, se establecen las siguientses conclusiones respecto al diseño estructural del edificio:

\begin{itemize}
    \item \textbf{Documento Bases de Cálculo:} Este informe constituye el marco normativo y técnico fundamental para el análisis del proyecto. Integra las disposiciones de las normas NCh433, NCh3171, NCh1537 y el Decreto Supremo Nº60[cite: 112, 114, 116, 119].
    
    \item \textbf{Periodo Aproximado Resultante:} El análisis dinámico preliminar arroja un periodo fundamental de vibración de $T = 1,05\text{ s}$. Este valor es consistente con la tipología de una estructura de 21 niveles basada en muros de corte de hormigón armado.
    
    \item \textbf{Influencia del Emplazamiento y Geotecnia:} La ubicación en Antofagasta (Zona Sísmica 3) impone la demanda sísmica más exigente, con una aceleración efectiva máxima de $A_{0} = 0,40g$. La clasificación preliminar como Suelo Tipo A favorece el diseño al minimizar la amplificación de ondas sísmicas.
    
    \item \textbf{Criterios de Uso de Hormigones:} Se define una estratificación de resistencias para optimizar la eficiencia estructural:
    \begin{itemize}
        \item \textbf{Subterráneo y niveles inferiores:} Hormigón de resistencia $f'_{c} = 35\text{ MPa}$ para absorber las máximas solicitaciones en la base.
        \item \textbf{Pisos superiores:} Hormigón de resistencia $f'_{c} = 25\text{ MPa}$ para reducir el peso propio en altura[cite: 80].
    \end{itemize}
\end{itemize}